% JETSPIN jet refinement documentation

\section{Dynamic refinement\label{sec:Dynamic-refinement}}

In the original Ren model, each bead represent a constant quantity
$V_{c}$ of jet volume, so that the relationship $\pi a_{i}^{2}\,l_{i}=\pi a_{0}^{2}l_{step}=V_{c}$
is valid for any $i-th$ bead. Despite such relationship can be sometimes
a useful advantage (as example, see Sec \ref{subsec:Introduction-of-dimensionless}),
the radius reduction $a_{i}$ of the electrospun nanofiber provides
a significant increase of the discretization length $l_{i}$ value
along the jet path (also greater than two order of magnitude). As
consequence, it was observed in literature \cite{kowalewski2009modeling}
that the jet discretization close the collector becomes rather coarse
to model efficiently the nanofiber. To get around this problem the
dynamic refinement procedure developed in Ref \cite{lauricella2016refinement}
can be used in JETSPIN in order to maintain the the length of the
elements $l_{i}$ below a prescribed characteristic length max $l_{max}$.
Whenever a bead length $l_{i}$ is larger than $l_{max}$, the nanofiber
is refined by discretizing uniformly the jet at the length step $l_{step}$
value given in input file. The discretization is performed by Akima
spline interpolations \cite{akima1970new,akima1991method} of the
main quantities describing the jet beads (positions, fiber radius,
stress, velocities). In order to perform the interpolation we proceed
following three steps. First of all, we introduce the variable $\lambda\in[0,1]$
to parametrize the jet path, where $\lambda=0$ identifies the nozzle,
and $\lambda=1$ the jet at the collector. For each bead we compute
its respective $\lambda_{i}$ value by using the formula

\begin{equation}
\lambda_{i}=\frac{1}{l_{path}}\sum_{k=1}^{i}l_{k}
\end{equation}
where $l_{path}=\sum_{k=1}^{N_{beads}}l_{k}$ is the jet path length
from the nozzle to the collector, and $i=1$ and $i=N_{beads}$ denote
the jet bead at the nozzle and collector, respectively ($i=0$ denotes
the nozzle). The set of ${\lambda_{i}}$ values represent the mesh
used to build the cubic spline. Given a generic quantity $y$, the
data $y_{i}=y\left(\lambda_{i}\right)$ are tabulated and used in
order to compute the coefficients of the \textit{Akima spline}, following
the algorithm reported in Ref \cite{akima1991method}. Then, we enforce
a uniform parametrization of the nanofiber by imposing all the lengths
of the elements equal to $l_{step}$. The new mesh ${\lambda_{i}^{*}}$
is defined as

\begin{equation}
\lambda_{0}^{*}=0\;\lambda_{i}^{*}=\frac{i}{N_{beads}^{*}},\;i=1,2,\ldots,N_{beads}^{*}\text{.}
\end{equation}
where $N_{beads}^{*}=l_{path}/l_{step}$ represents the number of
jet beads in the new representation. Finally, the new values $y\left(\lambda_{i}^{*}\right)$
are computed for any $i-th$ bead by the \textit{spline interpolation}.
The procedure is repeated for positions, fiber radius, stress and
velocities of the jet beads in order to provide the respective quantities
in the new mesh ${\lambda_{i}^{*}}$. It is worth to dressing that
the volume $V_{i}=\pi a_{i}^{2}\,l_{i}$ represented by each bead
is not anymore constant, and, consequently, $\pi a_{i}^{2}\,l_{i}\neq\pi a_{0}^{2}l_{step}$.
It is also possible keep few beads in their positions before and after
the refinement procedure in order to reduce possible errors introduced
by the interpolation procedure. In other words, these anchor beads
are manteined as knots of both the old and new meshes (see directive
\texttt{dynam refin anchor} in Tab \ref{tab:inputlist}). In the
following text we will refer to this emended version of the Ren model
as DyRen model. Further details on the implemented dynamic refinement
procedure are available in Ref \cite{lauricella2016refinement}.
